\chapter{Literature Review}

\section{Introduction}

The field of multi-agent systems and automated code generation has evolved significantly with the advent of Large Language Models (LLMs). This literature review examines the current state of the art in multi-agent frameworks, code generation techniques, and the challenges that our MetaGPT project aims to address.

Recent research has demonstrated remarkable progress in automated problem solving through societies of agents based on large language models. However, existing LLM-based multi-agent systems can only solve simple dialogue tasks effectively. Solutions to more complex tasks are complicated through logic inconsistencies due to cascading hallucinations caused by naively chaining LLMs.

The research landscape reveals several critical challenges that our project addresses:

\begin{itemize}
    \item \textbf{Logic Inconsistencies:} Existing systems suffer from cascading hallucinations when chaining LLMs, leading to unreliable outputs.
    \item \textbf{Lack of Structured Workflows:} Current multi-agent frameworks lack standardized operating procedures that guide agent interactions.
    \item \textbf{Communication Inefficiencies:} Agents often communicate through unconstrained natural language, leading to information distortion and inefficient collaboration.
    \item \textbf{Insufficient Self-Correction:} Most systems lack robust mechanisms for debugging and improving generated code.
\end{itemize}

Our MetaGPT framework directly addresses these challenges by incorporating Standardized Operating Procedures (SOPs) into multi-agent collaborations, enabling agents with human-like domain expertise to verify intermediate results and reduce errors.

\section{Key Papers Reviewed}

The literature on multi-agent systems and automated code generation provides a rich foundation for understanding the challenges and opportunities in this domain. Several key research papers have shaped the field of automated code generation and multi-agent systems, each addressing different aspects of the problem space.

\subsection{MetaGPT: Meta Programming for Multi-Agent Collaborative Framework (Hong et al., 2023)}
\textbf{Authors:} H. Hong, Y. Du, J. Pan, Z. Wu, and W. Chen

\textbf{Problem Addressed:} How can multiple LLM agents collaborate to solve complex software engineering tasks with reduced hallucination and higher consistency compared to naive multi-agent chat?

\textbf{Core Idea and Methodology:}
\begin{itemize}
    \item Proposes a ``meta programming'' view of multi-agent collaboration, where agents follow structured roles and produce explicit intermediate artifacts (e.g., requirement documents, design plans, and code deliverables) rather than relying only on free-form dialogue.
    \item Emphasizes \textbf{role specialization} (e.g., product-oriented analysis, architecture design, engineering implementation, and quality assurance) to mimic real software teams.
    \item Introduces \textbf{standardized workflows} to reduce information loss across agent handoffs and to make the process more auditable and repeatable.
\end{itemize}

\textbf{Strengths:}
\begin{itemize}
    \item Encourages structured communication through artifacts, improving traceability from requirements to implementation.
    \item Provides a clear decomposition of responsibilities across agents, making it easier to scale to multi-step software tasks.
    \item Highlights the importance of verification and quality checks as part of the workflow rather than an afterthought.
\end{itemize}

\textbf{Weaknesses / Limitations:}
\begin{itemize}
    \item Practical deployments still face engineering challenges such as persistence, observability, and tool/runtime integration for executing and validating generated projects.
    \item Performance and reliability depend on the quality of agent prompts, structured outputs, and the availability of mechanisms to recover from failures.
\end{itemize}

\textbf{Relevance to our project:} The proposed framework adopts this central idea by implementing a role-based pipeline (Manager, Architect, Engineer, QA) with explicit SOP definitions, graph-based orchestration, persistent storage of project state, and support services (e.g., sandboxed previews) to operationalize the research concepts into an end-to-end system.

\subsection{Evaluating Large Language Models Trained on Code (Codex, OpenAI 2021)}
\textbf{Authors:} Mark Chen, Jerry Tworek, Heewoo Jun, Qiming Yuan

\textbf{Problem Addressed:} Can large language models (trained on natural language) be adapted for code generation?

\textbf{Methodology:}
\begin{itemize}
    \item Fine-tuned GPT-3 on a massive GitHub code corpus to create Codex
    \item Generates code from natural language prompts and docstrings
    \item Benchmarked on HumanEval programming problems
\end{itemize}

\textbf{Strengths:}
\begin{itemize}
    \item Fluent code generation capabilities
    \item Effective for autocomplete, boilerplate, and simple scripts
    \item Demonstrated the potential of LLMs for code generation
\end{itemize}

\textbf{Weaknesses:}
\begin{itemize}
    \item Struggles with multi-file and system-level projects
    \item Prone to hallucination and insecure code generation
    \item Limited verification beyond sampling and reranking
\end{itemize}

\subsection{Competition-Level Code Generation with AlphaCode (DeepMind, 2022)}
\textbf{Authors:} Yujia Li, David Choi, Junyoung Chung, Nate Kushman, Julian Schrittwieser

\textbf{Problem Addressed:} Can LLMs solve competitive programming problems that require advanced reasoning, algorithm design, and coding?

\textbf{Methodology:}
\begin{itemize}
    \item Trained transformer-based models on competitive programming datasets
    \item Generated a very large pool of candidate programs through sampling
    \item Applied clustering and filtering to remove near-duplicates and select diverse solutions
    \item Used unit testing to automatically execute programs on hidden test cases
\end{itemize}

\textbf{Strengths:}
\begin{itemize}
    \item Strong performance on algorithmic challenges
    \item Effective verification via unit tests
    \item Robust solution space exploration through massive sampling
\end{itemize}

\textbf{Weaknesses:}
\begin{itemize}
    \item Extremely compute-intensive (generates hundreds of thousands of candidates)
    \item Focused on isolated problems, not full software projects
    \item Lacked structured planning or collaboration mechanisms
\end{itemize}

\subsection{Auto-GPT for Online Decision Making (2023)}
\textbf{Authors:} Hui Yang, Sifu Yue, Yunzhong He

\textbf{Problem Addressed:} Can LLMs operate autonomously, making decisions and using tools to complete multi-step goals without human step-by-step input?

\textbf{Methodology:}
\begin{itemize}
    \item Introduced Auto-GPT as an experimental open-source framework
    \item Uses chains of LLM prompts with a looping mechanism:
    \begin{itemize}
        \item Plan next step
        \item Execute (tool use, web browsing, file editing)
        \item Evaluate outcome
        \item Continue until goal achieved
    \end{itemize}
    \item Applied mainly to online decision-making tasks (web search, file operations, basic software tasks)
\end{itemize}

\textbf{Strengths:}
\begin{itemize}
    \item Handles long-horizon tasks beyond one-shot prompting
    \item Flexible and can adapt to new tasks via tool use
    \item Pioneered the autonomous LLM trend
\end{itemize}

\textbf{Weaknesses:}
\begin{itemize}
    \item Brittle and prone to getting stuck in loops
    \item No structured planning (SOPs) leading to ad-hoc decision-making
    \item Weak verification produces errors without correction
\end{itemize}

\subsection{MapCoder: Multi-Agent Code Generation for Competitive Problem Solving (2024)}
\textbf{Authors:} Yujia Li, David Choi, Junyoung Chung, Nate Kushman, Julian Schrittwieser

\textbf{Problem Addressed:} Can multi-agent LLMs collaborate to solve competitive programming problems efficiently?

\textbf{Methodology:}
\begin{itemize}
    \item Each agent specializes in a subtask:
    \begin{itemize}
        \item Planner $\rightarrow$ decomposes problem into smaller steps
        \item Solver agents $\rightarrow$ generate candidate code fragments
        \item Verifier $\rightarrow$ runs tests and identifies errors
    \end{itemize}
    \item Agents communicate iteratively, exchanging partial solutions
    \item Evaluated on competitive programming datasets (Codeforces, AtCoder)
\end{itemize}

\textbf{Strengths:}
\begin{itemize}
    \item Efficient multi-agent coordination reduces computational cost compared to AlphaCode's massive sampling
    \item Verification is embedded in the workflow
    \item Parallel agent execution allows scalability
\end{itemize}

\textbf{Weaknesses:}
\begin{itemize}
    \item Focused on competitive programming, not full software projects
    \item Limited grounding as agents rely mainly on LLM reasoning, not external documentation
\end{itemize}

\section{Literature Matrix}

Table \ref{tab:literature_matrix} provides a comprehensive comparison of the key papers reviewed, highlighting their methodologies, strengths, weaknesses, and relevance to our MetaGPT framework.

\begin{table}[H]
\centering
\caption{Literature Matrix: Comparative Analysis of Key Papers}
\label{tab:literature_matrix}
\resizebox{\textwidth}{!}{%
\footnotesize
\begin{tabular}{|l|p{2.2cm}|p{1.8cm}|p{2.4cm}|p{2.0cm}|p{2.2cm}|p{2.4cm}|}
\hline
\multicolumn{1}{|c|}{\textbf{Sl.No}} & \multicolumn{1}{c|}{\textbf{TITLE}} & \multicolumn{1}{c|}{\textbf{AUTHOR(S)}} & \multicolumn{1}{c|}{\textbf{METHODOLOGY}} & \multicolumn{1}{c|}{\textbf{DATASET/CONTEXT}} & \multicolumn{1}{c|}{\textbf{ADVANTAGES}} & \multicolumn{1}{c|}{\textbf{DISADVANTAGES}} \\
\hline
1 & Chain-of-Thought Prompting Elicits Reasoning in Large Language Models & Jason Wei et al. (2023) & Introduces chain-of-thought (CoT) prompting with intermediate reasoning steps & Arithmetic reasoning: GSM8K, SVAMP, ASDiv, AQUA, MAWPS; Commonsense reasoning: CSQA, StrategyQA, SayCan, BIG-Bench tasks & Enables multi-step reasoning in LMs (improves arithmetic, commonsense, symbolic tasks) & Only works well with very large models, High computational cost for inference with huge models \\
\hline
2 & ReAct: Synergizing Reasoning and Acting in Language Models & Shunyu Yao et al. (2023) & Introduces ReAct prompting: interleaves reasoning traces with actions & ALFWorld (text-based games), WebShop (online shopping with user instructions) & Combines internal reasoning with external actions, reducing hallucinations & More reasoning errors than CoT due to rigid structure (sometimes loops or repeats) \\
\hline
3 & CodeTree: Agent-guided Tree Search for Code Generation with Large Language Models & Jierui Li et al. (2025) & Framework for LLM agents using unified tree structure to explore search space for code generation & Function implementation tasks (HumanEval, MBPP, HumanEval+, MBPP+) and Program implementation tasks (CodeContests, APPS) & Flexible and efficient approach that avoids redundant exploration of solutions & Struggles with very difficult, competition-level problems \\
\hline
4 & Exploration of LLM Multi-Agent Application Implementation Based on LangGraph+Crew AI & Zhihua Duan and Jialin Wang (2024) & Use of LangGraph and CrewAI in multi agent system for code generation & Demonstrates framework capabilities through application cases in code generation, code review, ticket auditing and forwarding & Optimizes task execution and enhances system flexibility and scalability through real-time status data sharing and feedback mechanisms & The paper does not explicitly state any disadvantages of using the LangGraph and CrewAI frameworks together \\
\hline
5 & Enhancing Function-Calling Capabilities in LLMs: Strategies for Prompt Formats, Data Integration, and Multilingual Translation & Yi-Chang Chen et al. (2024) & Tuning-based approach to improve function-calling in Large Language Models & 110k instruction-following instances from Open ORCA (IF-110k) and 110k function-calling instances (FC-110k) from APIGen & Instruction-following data significantly improves function-calling accuracy and relevance detection & The Decision Token approach resulted in a slight decrease in function-calling accuracy \\
\hline
6 & Large Language Model based Multi-Agents: A Survey of Progress and Challenges & Taicheng Guo et al. (2024) & Summarizes commonly used datasets, benchmarks, and open-source frameworks & N/A & Highlights the advantages of LLM-based multi-agent systems over single-agent systems & Identifies the escalating complexity of managing a large number of agents as a critical problem \\
\hline
\end{tabular}%
}
\end{table}

\begin{table}[H]
\centering
\caption{Literature Matrix: Comparative Analysis of Key Papers (Continued)}
\label{tab:literature_matrix_continued}
\resizebox{\textwidth}{!}{%
\footnotesize
\begin{tabular}{|l|p{2.2cm}|p{1.8cm}|p{2.4cm}|p{2.0cm}|p{2.2cm}|p{2.4cm}|}
\hline
\multicolumn{1}{|c|}{\textbf{Sl.No}} & \multicolumn{1}{c|}{\textbf{TITLE}} & \multicolumn{1}{c|}{\textbf{AUTHOR(S)}} & \multicolumn{1}{c|}{\textbf{METHODOLOGY}} & \multicolumn{1}{c|}{\textbf{DATASET/CONTEXT}} & \multicolumn{1}{c|}{\textbf{ADVANTAGES}} & \multicolumn{1}{c|}{\textbf{DISADVANTAGES}} \\
\hline
7 & Agent AI with LangGraph: A Modular Framework for Enhancing Machine Translation Using Large Language Models & Jialin Wang and Zhihua Duan (2024) & Intent Agent to analyze input text and route to appropriate translation agent using LLMs & Uses English and French sentence pairs to train and test the model. A total of 75,000 words were selected for the experiment & High degree of flexibility, scalability, and accuracy in multilingual translation tasks & The model structure was relatively simple and the number of data sets was insufficient \\
\hline
8 & Retrieval Augmented Code Generation and Summarization & Md Rizwan Parvez et al. (2021) & Proposes two-step, retrieval-augmented framework called REDCODER with retriever and generator modules & CodeXGLUE and Concode benchmarks, framework was evaluated on code generation and summarization tasks in Java and Python & Enhances code generation and summarization by augmenting input with retrieved information & Computational expense of cross-encoder code retrievers and inability of sparse token-level retrievers (like BM25) to find semantically similar code and summaries \\
\hline
9 & Evaluating Large Language Models Trained on Code (Codex) & Mark Chen, Jerry Tworek et al. (2021) & Fine-tuned GPT-3 on a massive GitHub code corpus; benchmarked on HumanEval programming problems & HumanEval (164 hand-written programming problems) & Demonstrated the potential of LLMs for code generation; effective for autocomplete, boilerplate, and simple scripts & Struggles with multi-file and system-level projects; prone to hallucination and insecure code generation \\
\hline
10 & Competition-Level Code Generation with AlphaCode & Yujia Li, David Choi et al. (2022) & Trained transformer-based models on competitive programming datasets; generated large pool of candidates via sampling, then filtered using clustering and unit tests & CodeContests competitive programming dataset & Strong performance on algorithmic challenges; effective verification via unit tests; robust solution space exploration & Extremely compute-intensive; focused on isolated problems, not full software projects; no structured planning or collaboration \\
\hline
11 & Auto-GPT for Online Decision Making & Hui Yang, Sifu Yue, Yunzhong He (2023) & Chains LLM prompts in a plan-execute-evaluate loop with tool use (web browsing, file editing) & Online decision-making tasks including web search and file operations & Handles long-horizon tasks; flexible tool use; pioneered the autonomous LLM agent trend & Brittle and prone to getting stuck in loops; no structured SOPs; weak verification produces errors without correction \\
\hline
12 & MapCoder: Multi-Agent Code Generation for Competitive Problem Solving & Md. Ashraful Islam et al. (2024) & Multi-agent pipeline: Planner decomposes problem, Solver agents generate code, Verifier runs tests and feeds back iteratively & Competitive programming datasets: Codeforces, AtCoder, HumanEval, MBPP, APPS & Efficient multi-agent coordination; verification embedded in workflow; parallel execution allows scalability & Focused on competitive programming, not full software projects; limited grounding in external documentation \\
\hline
13 & MetaGPT: Meta Programming for Multi-Agent Collaborative Framework & H. Hong, Y. Du, J. Pan, Z. Wu, and W. Chen (2023) & Meta programming view: agents follow structured roles and produce explicit intermediate artifacts (requirements, design plans, code) rather than free-form dialogue & HumanEval, MBPP, and SoftwareDev software engineering benchmarks & Structured communication improves traceability; clear decomposition of responsibilities; verification and quality checks are part of the workflow & Practical deployments face challenges in persistence and observability; performance depends heavily on prompt quality and structured output reliability \\
\hline
\end{tabular}%
}
\end{table}

The literature matrix reveals several key insights:

\begin{itemize}
    \item \textbf{Evolution of Approaches:} From single LLM (Codex) to multi-agent systems (MapCoder), showing progression toward collaborative frameworks.
    \item \textbf{Verification Gaps:} Most approaches lack robust verification mechanisms, with only AlphaCode implementing comprehensive testing.
    \item \textbf{Structured Workflow Deficiency:} None of the existing approaches implement structured workflows or Standardized Operating Procedures (SOPs).
    \item \textbf{Context Awareness Limitations:} All approaches rely primarily on LLM reasoning without external knowledge integration.
    \item \textbf{Scalability Challenges:} Most approaches face scalability issues, either through computational intensity or lack of structured planning.
\end{itemize}

\section{MetaGPT Framework}

To address the limitations of existing multi-agent systems, particularly the logic inconsistencies and cascading hallucinations caused by naively chaining LLMs, we introduce the MetaGPT framework. Based on the foundational concepts proposed by Hong et al. (2023) in their paper "MetaGPT: Meta Programming for a Multi-Agent Collaborative Framework", MetaGPT incorporates efficient human workflows into LLM-based multi-agent collaborations through several key innovations:

\begin{itemize}
    \item \textbf{Standardized Operating Procedures (SOPs):} MetaGPT encodes human domain expertise and established SOPs into prompt sequences. This streamlines workflows and significantly improves the consistency and accuracy of task execution compared to free-form dialogue.
    \item \textbf{Role Specialization and Complete Workforce:} The framework utilizes an assembly line paradigm, assigning specialized roles such as Product Manager, Architect, Project Manager, Engineer, and QA Engineer. Each role follows specific constraints and responsibilities, mimicking a real-world software company to efficiently break down complex tasks.
    \item \textbf{Structured Communication Protocols:} Instead of unconstrained natural language chatter, agents communicate through structured artifacts like Product Requirement Documents (PRDs), system designs, APIs, and flowcharts. This minimizes ambiguities, reduces information distortion, and prevents the ``telephone game'' effect.
    \item \textbf{Publish-Subscribe Message Mechanism:} Agents share information through a global message pool. Rather than one-to-one communication, agents publish their structured outputs to the pool and subscribe to relevant messages based on their role profiles, avoiding information overload and improving efficiency.
    \item \textbf{Executable Feedback and Iterative Programming:} To ensure code quality and executability, MetaGPT employs a self-correction mechanism during runtime. Engineers generate code, write unit tests, and execute them. If errors occur, the agent debugs the code using historical execution memory and past artifacts, improving the code generation iteratively.
\end{itemize}

Extensive evaluations on collaborative software engineering benchmarks, such as HumanEval and MBPP, demonstrate that MetaGPT achieves state-of-the-art performance, validating the effectiveness of structured workflows and specialized multi-agent collaboration in complex automated programming tasks.

\section{Research Gaps and Our Contribution}

Grounded in the surveyed literature and current tool landscape, several concrete gaps can be identified that are directly relevant to our proposed project:

\begin{itemize}
    \item \textbf{Lack of end-to-end, production-oriented systems:} Many works focus either on solving isolated programming problems (e.g., competitive programming) or on early research prototypes of multi-agent systems. There is a shortage of production-ready platforms that take a natural language prompt all the way to a structured, full-stack project with APIs, UI, tests, and persistent storage.
    \item \textbf{Weak role modeling and orchestration:} Existing multi-agent approaches often describe roles conceptually, but do not provide a concrete, reusable orchestration layer where specialized agents are wired together with explicit state and transitions. Robust graph-based orchestration for software projects remains underexplored.
    \item \textbf{Limited observability and traceability:} Most systems do not expose detailed reasoning traces, intermediate artifacts, or execution timelines to end users. This makes it difficult for developers to understand why a system produced a particular codebase or to debug agent behavior.
    \item \textbf{Insufficient integration with modern web development stacks:} Few research systems are tightly integrated with contemporary stacks such as FastAPI on the backend and Next.js/React on the frontend, with opinionated project structures that align with how teams actually build and ship applications.
    \item \textbf{Underutilization of retrieval for code-level context:} While retrieval-augmented generation has been explored for natural language and small code snippets, there is limited work on integrating RAG deeply into agentic workflows for managing entire projects, file trees, and prior artifacts.
\end{itemize}

These gaps motivate the design choices in our framework, which explicitly targets end-to-end, observable, and stack-aware agentic code generation.